Resuelve los siguientes problemas:
        \begin{multicols}{2}
            \begin{parts}
                \part Una casa tiene una alberca circular de 6 metros de diámetro. Calcula el área de la alberca.
                \fillin[$28.26$ m$^2$][0in]
                
				\begin{solutionbox}{1cm}
					$A=\pi r^2 = \pi (3)^2 = 28.26$ m$^2$
				\end{solutionbox}

				\part El radio de una rueda es de 32 centímetros, ¿cuántos centímetros habrá recorrido esa rueda después de haber dado 22 vueltas?
                \fillin[$70737.92$ cm][0in]

				\begin{solutionbox}{1.2cm}
					$C=2\pi r = 2\pi (32) = 201.06$ cm\\
					$22(201.06) = 70737.92$ cm
				\end{solutionbox}

                \columnbreak

                \part Calcula el área de un parque que tiene un radio de 170 metros.
                \fillin[$90746$ m][0in]
                
				\begin{solutionbox}{1cm}
					$A=\pi r^2 = \pi (170)^2 = 90746$ m$^2$
				\end{solutionbox}

				\part Daniel tiene un terreno circular con un radio de 6 metros al cual le desea poner una barda en su periferia, si el precio por metro de barda es de 124 pesos. ¿Cuánto pagará en total por poner la barda?
                \fillin[$\$4,672.32$ pesos][0in]

				\begin{solutionbox}{1.2cm}
					$P=2\pi r = 2\pi (6) = 37.68$ m\\
					$37.68(124) = \$4672.32$ pesos
				\end{solutionbox}
            \end{parts}
        \end{multicols}